\documentclass{standalone}
\usepackage{tikz}
\usetikzlibrary{intersections, calc}
\usepackage{hsbcommon}
    \begin{document}
\begin{tikzpicture}
        \draw[->] (0,0) node[below] {\phantom{$X_i$}} -- (8,0) node[right]  {B};
        \draw[->] (0,0) -- (0,8) node[above]  {A};

        \def\mya{A}
        \def\myb{B}
        \def\mys{S}
        \def\myp{P}
        \def\myq{Q}

        \foreach [count=\i] \x/\y/\mycolor in {7/7/red%
            ,5/5/blue%
            }
        {       
                \begin{scope}[color=\mycolor]
                    \coordinate (\mya) at (0,\y);
                    \coordinate (\myb) at (\x,0);
                    % \node[above right] at (\myb) {$B_\i$};
             
                    % \draw [name path=budgetgerade] plot [smooth] coordinates {(\mya) (\myb) };
            
                    \coordinate (\mys) at ($(\mya)!0.4!(\myb)$);  %% punkt (fast) in der mitte
                    % \node[above right] at (\mys) {$S_\i$};
                    
                    % \draw [dashed]  (\mys) -- (\mys |-,0);
                    % \draw [dashed]  (\mys) -- (0, |- \mys);
                    
                    % \filldraw (\mys) circle (2pt);
                    \coordinate (\myp) at ([shift=({-2cm,3.5cm})]\mys); %% punkt links oben
                    \coordinate (\myq) at ([shift=({3.5cm,-2cm})]\mys); %% punkt rechts unten
                    
                    \node[right] at (\myq) {$U_\i$};
            
                    \draw [name path=indifferenzkurve] plot [smooth, tension=1] coordinates {(\myp) (\mys) (\myq) };
                \end{scope}
        }        
        
    \end{tikzpicture}
\end{document}
